\documentclass[11pt,a4paper]{article}
\usepackage[T1]{fontenc}
\usepackage{lmodern,microtype}
\usepackage[margin=25mm]{geometry}
\usepackage{amsmath,amssymb,amsthm,mathtools}
\usepackage{booktabs,tabularx,array}
\usepackage{enumitem}
\usepackage[numbers,sort&compress]{natbib}
\usepackage{xcolor}
\usepackage{xurl}
\usepackage[colorlinks=true,linkcolor=blue!45!black,citecolor=blue!45!black,urlcolor=blue!45!black]{hyperref}
\usepackage{fancyhdr}
\pagestyle{fancy}
\fancyhf{}
\fancyhead[L]{\small Supplement S1: weighted GRIP complexity}
\fancyfoot[C]{\thepage}
\setlength{\headheight}{14pt}
\setlength{\emergencystretch}{2em}
\setlist{itemsep=3pt,topsep=5pt}
\newcommand{\bigO}{\mathcal O}
\newcommand{\ball}{\mathcal B}
\newcommand{\work}{\mathcal W}
\newcommand{\code}[1]{\texttt{#1}}
\renewcommand{\thesection}{S1.\arabic{section}}
\renewcommand{\theequation}{S1.\arabic{equation}}
\renewcommand{\thepage}{S1--\arabic{page}}
\title{Supplement S1\\Search-sensitive complexity of weighted GRIP}
\author{Pawel Gajer \and Jacques Ravel}
\date{}
\begin{document}
\maketitle

\section{Scope of the result}

Weighted GRIP uses two forms of local shortest-path search. Filtration
searches stop at a prescribed weighted radius, and insertion searches stop
after their neighborhood and placement-anchor requirements are met. Charging
a full-graph traversal to every query therefore gives an upper bound that
does not reflect these stopping rules. This supplement derives an estimate
in terms of the search work actually performed.

The analysis concerns the 2D/3D compiled weighted core in \texttt{grip} 0.2.0
\citep{gripcode2026}. For filtration, the estimate is controlled by the overlap
of equal-radius balls around the selected centers. For insertion, it depends
on the number of vertices traversed before enough retained neighbors and
anchors are found. Under bounded degree, bounded filtration-ball overlap,
suitable coarse-vertex coverage, logarithmic hierarchy depth, and fixed
tuning parameters, these estimates give a conditional
$\bigO(n\log^2 n)$ expected-time upper bound, including filtration and
coarsest initialization. The expectation concerns sampled repulsion.
Geometric contraction of level sizes additionally gives linear peak storage.

The search-sensitive accounting follows the implementation. The favorable
asymptotic bound requires the additional geometric assumptions stated in
Section~\ref{sec:conditional}; it is not an unconditional guarantee for
weighted graphs. No new scalability experiment is reported.

\section{Scope and the cost of one search}

The input is a connected, undirected, simple graph $G=(V,E,\ell)$ with $n=|V|$
vertices, $M=|E|$ edges, and strictly positive edge lengths. Distances below are
measured after any global edge-length normalization. The analysis starts with
the normalized adjacency representation supplied to the compiled core; R input
validation and normalization are outside its scope. Dimension is fixed at two
or three. The default first-anchor rule and final local Fruchterman--Reingold
pass are used; optional LGKK, post-layout edge-KK, trace storage, and alternative
anchor-selection strategies are excluded.

A search $s$ examines $a_s$ adjacency incidences, performs $p_s$ heap pushes and
pops in total, and reaches a maximum heap size $H_s$. The binary heap includes
stale entries left by improved tentative distances. Its traversal work is
bounded by
\begin{equation}
 \tau_s=\bigO\!\left(1+a_s+p_s\log(2+H_s)\right).
 \label{eq:query}
\end{equation}
Each successful relaxation causes at most one push, and every pop removes a
previously pushed item. Consequently $p_s\leq 2(a_s+1)$ and a convenient upper
bound is $\tau_s=\bigO((1+a_s)\log(2+H_s))$. This counts boundary edges as well
as internal edges: an edge examined at the search frontier costs work even when
its other endpoint is never settled. Destroying an unfinished heap is also
covered by the number of pushes. Visitor work for filling caches is accounted
for separately below.

\clearpage
\section{Filtration: radius-limited work and ball overlap}

\subsection{The actual construction searches}

The weighted construction begins with $C_0=V$. At scale $i$, it greedily selects
centers $C_i\subseteq C_{i-1}$ and removes candidates at distance at most
$r_i=2^{i-1}r_0$ from each newly selected center; the implementation uses
$r_0=1$ in normalized units. Its traversal runs in the original graph, not in
the subgraph induced by the current candidates. Every center therefore incurs
one search of the ball $\ball(v,r_i)=\{x:d_G(v,x)\leq r_i\}$. The selected centers
are mutually more than $r_i$ apart in the ideal shortest-path model.

The $t$ nonempty construction scales are kept distinct from the final layout
hierarchy. The code can discard an overcoarse level and append a fixed-size
initialization subset; that appended subset is not necessarily separated at
the next geometric radius. The packing arguments here apply to the genuine
greedy center sets $C_i$, not to that synthetic last subset.

The graph-volume of a vertex set $A$ is
$w(A)=\sum_{x\in A}(1+\deg_G(x))$. It counts vertices and the adjacency
incidences that must be scanned. At level $i$, the total searched volume is
\begin{equation}
 W_i=\sum_{v\in C_i} w\bigl(\ball(v,r_i)\bigr).
 \label{eq:volume}
\end{equation}
With $H_i$ the largest heap during those searches, the radius-sensitive bound is
\begin{equation}
 T_F=\bigO\!\left(n+\sum_{i=1}^{t} W_i\log(2+H_i)\right).
 \label{eq:filtration}
\end{equation}
It does not replace a small radius-ball with the entire graph. The reusable
distance and epoch arrays cost $\bigO(n)$ initially, not $\bigO(n)$ per query.

\subsection{Why overlap is the decisive quantity}

A vertex $x$ is reached by
$\mu_i(x)=|\{v\in C_i:d_G(v,x)\leq r_i\}|$ center searches at scale $i$.
Reversing the order of summation gives the exact volume identity
\begin{equation}
 W_i=\sum_{x\in V}(1+\deg_G(x))\mu_i(x)
     =(n+2M)\,\overline\mu_i,
 \label{eq:overlap}
\end{equation}
where $\overline\mu_i$ is the degree-weighted mean overlap. Thus
\begin{equation}
 \boxed{T_F=\bigO\!\left(n+(n+2M)
       \sum_{i=1}^{t}\overline\mu_i\log(2+H_i)\right).}
 \label{eq:overlap-bound}
\end{equation}
Even a large ball is affordable when few centers generate overlapping searches.
Small balls and decreasing center counts should not be analyzed independently.

For a sufficient geometric condition, suppose every graph-metric ball of radius
$r$ is covered by at most $\lambda$ balls of radius $r/2$, uniformly over the
relevant scales. This is the doubling condition, with doubling constant
$\lambda$. The centers counted by $\mu_i(x)$ lie in $\ball(x,r_i)$ and are
pairwise more than $r_i$ apart. No two can lie in the same radius-$r_i/2$ cover
ball, by the triangle inequality. Therefore $\mu_i(x)\leq\lambda$, yielding
$T_F=\bigO(n+\lambda t(n+M)\log(2+M))$. Bounded graph degree alone does not
imply this bounded-overlap condition.

\clearpage
\section{Insertion: how far must the search go?}

\subsection{Neighbors at every relevant level, plus anchors}

The final layout hierarchy is $V_0=V\supseteq V_1\supseteq\cdots\supseteq V_h$,
with level sizes $N_j=|V_j|$. The top set has $q=N_h$ initialization vertices.
A vertex introduced at level $d(v)$ belongs to every $V_j$ for $j\leq d(v)$.
The retained-neighbor budget at level $j$ is $b_j=\min(b,N_j-1)$, where
$b$ is \code{num\_nbrs}. Insertion searches simultaneously fill all these caches
and locate $a$ anchors from higher levels. Under the default
\code{any\_higher} rule, eligible anchors are precisely $V_{d(v)+1}$.

The code's stopping condition is slightly stricter than ``the caches contain
$b_j$ entries.'' A full cache is marked complete only on the next eligible
vertex encountered. For level $j$, a sufficient stopping distance is therefore
the distance $\eta_j(v)$ to the $(b_j+1)$st other vertex in $V_j$. If there is no
such vertex, $\eta_j(v)=\infty$. The distance to the $a$th eligible anchor is
$\alpha(v)$, also infinite if too few exist. Under the default first-anchor
rule, a sufficient radius for completion is
\begin{equation}
 R_v=\max\!\left\{\alpha(v),\max_{0\leq j\leq d(v)}\eta_j(v)\right\}.
 \label{eq:stop-radius}
\end{equation}
The actual search can stop within the final distance tie, so the whole ball is
an upper envelope, not a requirement to visit every point at that distance.
If a completion condition is impossible, the finite graph is exhausted.

Searching for $b$ retained neighbors does \emph{not} imply that only $b$ graph
vertices are visited. Many nearby vertices may be absent from a coarse level.
Likewise, finding anchors can require searching farther than filling a local
cache. Equation~\eqref{eq:stop-radius} includes both constraints.

\subsection{The prefix-sensitive bound}

The settled prefix $P_v$ contains the vertices processed through termination,
including the root. When a callback or cap terminates the search, its last
vertex is not expanded. If $s_v=|P_v|$ and the maximum degree is $\Delta$,
then $a_v\leq\Delta s_v$ and $H_v\leq 1+\Delta s_v$. Hence the insertion-search
work, excluding cache writes and placement, satisfies
\begin{equation}
 T_N^{\rm search}=\bigO\!\left(
   \sum_{v\notin V_h}(1+\Delta s_v)
        \log(2+\Delta s_v)\right).
 \label{eq:prefix}
\end{equation}
For nonuniform degrees, use the actual $a_v,p_v,H_v$ in
Equation~\eqref{eq:query}. A heap can contain endpoints outside $P_v$; the bound
charges their insertion to scanned frontier edges, rather than assuming the
heap contains only settled vertices.

\subsection{An explicit approximate search cap}

When \code{metric\_neighbor\_cap} is a positive $C$, insertion stops after at
most $C$ non-root settlements, even if caches or anchors are incomplete. Thus
$s_v\leq C+1$ and Equation~\eqref{eq:prefix} becomes
$\bigO((n-q)(1+\Delta(C+1))\log(2+\Delta(C+1)))$ for traversal work. This is
an approximation mode. It does not limit filtration searches or the initial
full searches, and it does not bound scanned edges independently of degree.

\clearpage
\section{Complete core cost and memory}

\subsection{The full searches that remain}

For each of the $q$ top-level vertices, \code{metric\_me\_init\_v1()} performs
an uncapped, full-graph traversal, retaining selected neighbors while also
collecting the maximum distance used to size the initial layout box. Its visitor
always returns false. These searches contribute
\begin{equation}
 T_I^{\rm search}=\bigO\!\left(q(n+M)\log(2+M)\right)
 \label{eq:initial}
\end{equation}
as a general binary-heap upper bound. Their actual cost can be smaller when
the heap frontier stays small. The default top-level budget is 24 vertices, but a user-chosen
$q$ that increases with $n$ must remain in the bound.

\subsection{Cache work is amortized, not multiplied by every settlement}

Cache-vector headers number
$U=\sum_{v\in V}(d(v)+1)=\sum_{j=0}^{h}N_j$. The maximum number of stored
neighbor entries is
\begin{equation}
 K=\sum_{v\in V}\sum_{j=0}^{d(v)}b_j
   =\sum_{j=0}^{h}N_jb_j\leq bU.
 \label{eq:cache}
\end{equation}
Each iteration of the callback's inner loop either appends one retained entry
or advances its monotone lowest-unfinished-level index. All callbacks together
therefore add $\bigO(K+U+\sum_s |P_s|)$ work. There is no need to charge every
settled vertex an extra factor $h$. The last term is already covered by the
heap-pop counts in Equation~\eqref{eq:query}.

\subsection{Placement, refinement, and the aggregate bound}

With a fixed first-anchor count and a fixed number of local placement steps,
$T_{\rm place}=\bigO(n)$. At level $j$, the algorithm performs $R_j$ refinement
rounds. With fixed repulsion sample size and fixed exact-repulsion threshold,
the expected refinement cost is
\begin{equation}
 \mathbb E T_{\rm ref}=\bigO\!\left(
   R_0M+\sum_{j=0}^{h}R_jN_j(b_j+1)\right).
 \label{eq:refine}
\end{equation}
The expectation refers to the usual uniform-draw model for rejection sampling
of distinct repulsion partners. For a deterministic operation count, record
the actual number of attempted draws instead. Unbounded exact global repulsion
or sample sizes increasing with $n$ are outside this formula.

Across filtration, initialization, and insertion, the search calls form a set
$\mathcal Q$. Their realized counts give the principal implementation estimate:
\begin{equation}
 \boxed{T_{\rm core}=\bigO\!\left(
 n+M+\sum_{s\in\mathcal Q}
 [1+a_s+p_s\log(2+H_s)]+K+U+T_{\rm place}+T_{\rm ref}\right).}
 \label{eq:total}
\end{equation}
The notation $\mathcal Q$ indexes calls, so the same root can occur several
times. The graph, scratch arrays, cache headers, retained entries, and largest
live heap require $\bigO(n+M+U+K+\max_s H_s)$ peak space. For a simple graph
under standard Dijkstra accounting, $H_s=\bigO(M+1)$, reducing this to
$\bigO(n+M+U+K)$. If $N_{j+1}\leq\theta N_j$ for a constant $\theta<1$, then
$U=\bigO(n)$ and $K=\bigO(bn)$: fixed-$b$ cache storage is linear, not
necessarily $n\log n$.

\clearpage
\section{A conditional near-linear estimate}\label{sec:conditional}

\subsection{Additional assumptions about coarse-vertex coverage}

Bounded overlap controls filtration but does not, by itself, guarantee that
every insertion query finds enough coarse neighbors nearby. A sufficient
coverage condition is that a vertex inserted at level $j<h$ settles at most
\begin{equation}
 s_v\leq c\!\left[
   (b+1)\frac{n}{N_j}
   +a\frac{n}{N_{j+1}}+1\right],
 \label{eq:coverage}
\end{equation}
for a constant $c$ independent of $n$ and the level. The two fractions express
the amount of original-graph material searched per retained neighbor and per
higher-level anchor. This is an explicit assumption about actual search
prefixes, motivated by regularly sampled geometries and well-distributed nets;
it is not a consequence of bounded degree, nor a proved property of every
weighted GRIP hierarchy. It must include the appended top-level subset and
large distance ties if they affect stopping.

Suppose also that adjacent level sizes satisfy $N_j/N_{j+1}\leq\gamma$ for a
constant $\gamma$. Since $N_j-N_{j+1}$ vertices are inserted at level $j$,
Equation~\eqref{eq:coverage} gives
\begin{align}
 \sum_{v\notin V_h}s_v
 &\leq c' (b+1+a\gamma)n
     \sum_{j=0}^{h-1}\frac{N_j-N_{j+1}}{N_j} \notag\\
 &\leq c'(b+1+a\gamma)nh .
 \label{eq:coverage-sum}
\end{align}
The constant $c'$ absorbs the additive one. Combining this with bounded degree
in Equation~\eqref{eq:prefix} yields
$T_N^{\rm search}=\bigO((b+1+a\gamma)nh\log n)$.

\subsection{Putting the assumptions together}

For a bounded-degree graph, bounded filtration overlap $\lambda$, coverage
as in Equation~\eqref{eq:coverage}, and fixed tuning parameters, the search work
is bounded by
\begin{equation}
 \bigO\!\left(n[\lambda t+(b+1+a\gamma)h+q]\log n\right).
 \label{eq:conditional}
\end{equation}
If $t,h=\bigO(\log n)$ and the constants do not grow with $n$, the full core is
$\bigO(n\log^2 n)$ in expectation for the sampled-repulsion part. Filtration
and uncapped initialization are included. Under geometric contraction of the
level sizes, peak space is $\bigO(n)$ for fixed $b$, and refinement with fixed
round budgets is also $\bigO(n)$ in expectation.

This estimate explains why radius limits and early stopping can make weighted
GRIP scalable without an all-pairs distance matrix. It does not assert that
all weighted graphs satisfy its geometric assumptions, or that the displayed
upper bound is tight. With smaller actual heaps or search volumes, the bound
in Equation~\eqref{eq:total} is sharper still.

\paragraph{A one-dimensional example.}
For a uniformly weighted path with regularly spaced nested nets,
$|C_i|=\bigO(n/r_i)$ and each radius-ball contains $\bigO(r_i)$ vertices. The
heap frontier has at most two entries. Filtration therefore costs $\bigO(n)$
per scale, or $\bigO(n\log n)$ in total, without a heap logarithm. If the
retained levels have comparable spacing and bounded size ratios, insertion
searches likewise cost $\bigO(n)$ per scale for fixed $a,b$. The fixed number
of full initialization traversals costs $\bigO(n)$. The resulting
$\bigO(n\log n)$ estimate illustrates that radius-limited searches need not
be constant-size to produce near-linear overall work.

\clearpage
\section{Qualifications and comparison with the original bound}

\subsection{Weighted scale depth is not automatically logarithmic in \texorpdfstring{$n$}{n}}

For weighted diameter $D_w$, doubling radii suffice after
$t=\bigO(1+\log_+(D_w/r_0))$ construction scales, where
$\log_+x=\max(0,\log_2x)$. With arbitrary positive edge lengths, $D_w/r_0$
need not be polynomial in $n$. Median normalization does not prevent a very
large edge-length ratio. The present 2D/3D implementation allocates only
$\lfloor\log_2n\rfloor+2$ level slots and throws an error if the construction
needs more. Thus successful runs have logarithmic allocated depth; that cap is
not a theorem that every weighted input will successfully produce such a
hierarchy. The conditional estimate above assumes successful construction and
the stated depth bounds.

\subsection{Numerical and optional-mode qualifications}

Equation~\eqref{eq:total} is an operation-count estimate for the calls performed.
The metric-ball and packing derivations additionally use the ideal positive-
weight shortest-path model. The C++ code uses absolute and relative comparison
tolerances of order $10^{-10}$, so exact-radius identities should be understood
away from numerically unresolved boundary cases. If those cases matter, use
actual examined sets and multiplicities, or inflate radii and recheck the
separation-to-radius ratio; do not assume exact separation from floating-point
comparisons alone.

Distance scratch arrays are tagged by a 32-bit epoch. A wraparound incurs
an additional $\bigO(n)$ reset per $2^{32}$ traversal calls. A literal
machine-level count can add $\bigO(n\lfloor|\mathcal Q|/2^{32}\rfloor)$ to
Equation~\eqref{eq:total}. The usual word-RAM analysis treats the epoch as large
enough for the run.

The default first-anchor strategy uses a fixed number of anchors. Distance-band
strategies may retain large tied candidate sets and perform additional subset
selection. They need their own placement terms. Optional LGKK, later edge-KK,
exact all-vertex repulsion, and retained trace frames likewise require separate
costs. These qualifications prevent a bound for the base initializer from being
misread as a bound for every exported workflow.

\subsection{What follows, and what does not}

The published original GRIP result is $\Theta(nk^2)$ time and $\Theta(nk)$
space for bounded-degree graphs after filtration construction, with
$k=\log\delta(G)$ for an MIS filtration, where $\delta(G)$ is the unweighted
graph diameter \citep{gajer2004}. Its scale-dependent neighbor budgets differ
from the fixed per-cache cap $b$ in the current implementation; the historical
space bound therefore need not coincide with Equation~\eqref{eq:cache}.
The earlier GRIP
paper already describes filtration cost through the sizes of its partial BFS
trees, rather than only through a full-graph worst case \citep{gajer2002}.
The overlap and early-stopping estimates in this supplement extend that way of
accounting to the weighted implementation; they are derivations in this supplement, not
claims attributed to those papers.

Without locality assumptions, replacing every search volume by the whole graph
recovers the loose envelope
$\bigO(n(t+1)(n+M)\log n)$ for the search work. With $t=\bigO(\log n)$ it becomes
a $\bigO(n(n+M)\log^2 n)$ estimate. That upper bound is neither a
prediction of typical runtime nor a lower bound showing that weighted GRIP
must be quadratic. Conversely, the more favorable conditional result does not
establish that weighted GRIP is always faster than metric MDS. The defensible
conclusion is that its cost is governed by explored weighted neighborhoods,
their overlap across centers, and the density of retained coarse vertices.

\clearpage
\section{Implementation correspondence}

The operation counts above can be evaluated by recording, for each search,
its phase, root, filtration level, radius, adjacency scans $a_s$, heap
operations $p_s$, maximum heap occupancy $H_s$, and stopping reason.
Insertion-prefix sizes and cache and anchor completion states distinguish
early termination from full-graph exhaustion. Level-specific volumes $W_i$
and overlap $\overline\mu_i$ measure whether coarsening offsets increasing
search radius. Such measurements would test the relevance of the assumptions
to a given graph family; the conditional analysis alone does not establish
an observed scaling law.

The source locations below identify the operations analyzed. The inspected
snapshot, revision \code{cebcc62}, has the same
\nolinkurl{src/MishWeighted.cpp}, \nolinkurl{src/DrawGraph.cpp},
\nolinkurl{src/MishSupport.cpp}, \nolinkurl{src/grip_rcpp.cpp}, and
\nolinkurl{R/grip_layout_weighted.R} files as release \code{v0.2.0}
\citep{gripcode2026}. The bibliography links the full immutable commit
identifier; the correspondence is based on source inspection, not timing data.

\begin{center}
\small
\begin{tabularx}{\textwidth}{@{}p{0.40\textwidth}X@{}}
\toprule
Source location & Relevant behavior \\
\midrule
\nolinkurl{src/MishWeighted.cpp}, lines 28--115
& Reused epoch arrays, binary heap, cutoff, visitor stop, and optional settlement cap. \\
\addlinespace
Same file, lines 190--282
& One radius query per selected center; uncapped full query for each initialization root. \\
\addlinespace
Same file, lines 726--874
& Cache filling, extra eligible-vertex completion check, anchor selection, and early termination. \\
\addlinespace
\nolinkurl{src/DrawGraph.cpp}, lines 187--256
& Logarithmic allocation, per-level neighbor budgets, and the initialization-root loop. \\
\addlinespace
\nolinkurl{src/MishSupport.cpp}, lines 658--749
& Exact/sampled repulsion switch and rejection sampling of distinct partners. \\
\addlinespace
\nolinkurl{src/MishWeighted.cpp}, lines 1063--1369
& Final edge attraction, cached repulsion, per-level rounds, and once-only insertion. \\
\bottomrule
\end{tabularx}
\end{center}

\bibliographystyle{plainnat}
\bibliography{S1-weighted-grip-complexity}
\end{document}
